\section{Conclusion}

Boundary Density Likelihood addresses a common but consequential mismatch in time-series detection: models are often trained as samplewise segmenters, but evaluated by the timing of ranked detections.
It changes the unit of supervision from state occupancy to expected scored events on the model output timeline, then fits the corresponding conditional mean with a Poisson log score.
Across the sleep and CHB-MIT benchmarks, changing this training signal improves event AP under matched architectures, scorers, and post-processing.
On sleep, the gain is largest at strict matching tolerances, persists across multiple backbones, and is not reproduced by segmentation reweighting or smoothing choices alone.

The significance is practical as much as mathematical.
In health recordings, the few scored times often determine which episode is reviewed, which alert is raised, or which boundary enters a clinical workflow.
A detector that is good at labeling the interior of an interval can still be costly if it ranks the boundary late, duplicates it, or misses a short event.
BDL makes the sequence output estimate the temporal detections that event-level metrics and downstream review systems consume, while retaining the computational convenience of dense sequence prediction.
The state-reconstruction diagnostic keeps the scope clear: when the deployed target is samplewise occupancy, state labels remain the direct supervision.
When the deployed target is a ranked set of event times, the training signal should preserve and estimate those events.
